\documentclass[11pt,a4paper]{article}

\usepackage[utf8]{inputenc}
\usepackage[T1]{fontenc}
\usepackage[spanish]{babel}
\usepackage{amsmath,amssymb,mathtools}
\usepackage{geometry}
\usepackage{enumitem}
\usepackage{array}
\usepackage{lmodern}
\usepackage{microtype}
\usepackage{xcolor}

\geometry{margin=2.4cm}
\setlength{\parindent}{0pt}
\setlength{\parskip}{0.5em}

\newcommand{\R}{\mathbb{R}}
\newcommand{\K}{\mathbb{K}}
\newcommand{\Pn}{\mathcal{P}}
\newcommand{\Span}{\operatorname{L}}
\newcommand{\rg}{\operatorname{rg}}
\newcommand{\dimn}{\operatorname{dim}}

\begin{document}

\begin{center}
	{\LARGE\bfseries Álgebra Lineal}\\[0.4em]
	{\Large Tema 1: Espacios Vectoriales --- Soluciones}
\end{center}

\section{Ejercicio 1}
	Sea:

	\[
	S=
	\left\{
	(x,y,z)\in\R^3:
	2x-y+3z=0
	\right\}
	\]

	\subsection{Solución}

	El vector nulo pertenece a $S$:

	\[
	(0,0,0)\in S
	\]

	Sean $u=(x,y,z)\in S$ y $v=(x',y',z')\in S$. Entonces:

	\[
	2x-y+3z=0
	\]

	\[
	2x'-y'+3z'=0
	\]

	Para $\alpha,\beta\in\R$:

	\[
	2(\alpha x+\beta x')
	-(\alpha y+\beta y')
	+3(\alpha z+\beta z')
	\]

	\[
	=
	\alpha(2x-y+3z)
	+
	\beta(2x'-y'+3z')
	=
	0
	\]

	Por tanto:

	\[
	\boxed{S\text{ es subespacio vectorial de }\R^3}
	\]

	Despejando:

	\[
	y=2x+3z
	\]

	\[
	(x,y,z)
	=
	x(1,2,0)+z(0,3,1)
	\]

	\[
	\boxed{
	S=
	\Span\{(1,2,0),(0,3,1)\}
	}
	\]

\section{Ejercicio 2}
	Sea:

	\[
	S=
	\left\{
	(x,y)\in\R^2:
	xy=0
	\right\}
	\]

	\subsection{Solución}

	Aunque:

	\[
	(0,0)\in S
	\]

	el conjunto no es cerrado para la suma. Por ejemplo:

	\[
	(1,0)\in S,
	\qquad
	(0,1)\in S
	\]

	pero:

	\[
	(1,0)+(0,1)
	=
	(1,1)
	\notin S
	\]

	porque:

	\[
	1\cdot1\neq0
	\]

	Por tanto:

	\[
	\boxed{
	S\text{ no es subespacio vectorial de }\R^2
	}
	\]

\section{Ejercicio 3}
	Sea:

	\[
	S=
	\left\{
	A\in M_{2\times2}(\R):
	\det A=0
	\right\}
	\]

	\subsection{Solución}

	La matriz nula pertenece a $S$.

	Sin embargo, consideramos:

	\[
	A=
	\begin{pmatrix}
	1&0\\
	0&0
	\end{pmatrix},
	\qquad
	B=
	\begin{pmatrix}
	0&0\\
	0&1
	\end{pmatrix}
	\]

	Se tiene:

	\[
	\det A=\det B=0
	\]

	pero:

	\[
	A+B=
	\begin{pmatrix}
	1&0\\
	0&1
	\end{pmatrix}
	\]

	y:

	\[
	\det(A+B)=1
	\]

	Por tanto:

	\[
	\boxed{
	S\text{ no es subespacio vectorial}
	}
	\]

\section{Ejercicio 4}
	Determinar $x$ para que:

	\[
	(1,x,5)
	\in
	\Span\{(1,2,3),(1,1,1)\}
	\]

	\subsection{Solución}

	Deben existir $\alpha,\beta\in\R$ tales que:

	\[
	(1,x,5)
	=
	\alpha(1,2,3)
	+
	\beta(1,1,1)
	\]

	Por coordenadas:

	\[
	\alpha+\beta=1
	\]

	\[
	3\alpha+\beta=5
	\]

	Restando:

	\[
	2\alpha=4
	\]

	\[
	\alpha=2,
	\qquad
	\beta=-1
	\]

	Entonces:

	\[
	x=2\alpha+\beta
	=
	4-1
	=
	3
	\]

	\[
	\boxed{x=3}
	\]

\section{Ejercicio 5}
	Determinar si:

	\[
	A=
	\{p\in\Pn_2(x):p(0)=0,\ p'(0)=0\}
	\]

	y:

	\[
	B=
	\{p\in\Pn_2(x):p(0)=0,\ p'(0)=1\}
	\]

	son subespacios.

	\subsection{Solución}

	Sea:

	\[
	p(x)=a+bx+cx^2
	\]

	Para $A$:

	\[
	p(0)=a=0
	\]

	\[
	p'(x)=b+2cx
	\]

	\[
	p'(0)=b=0
	\]

	Por tanto:

	\[
	A=
	\{cx^2:c\in\R\}
	=
	\Span\{x^2\}
	\]

	\[
	\boxed{
	A\text{ es subespacio vectorial}
	}
	\]

	Para $B$:

	\[
	p(0)=0
	\quad\Longrightarrow\quad
	a=0
	\]

	\[
	p'(0)=1
	\quad\Longrightarrow\quad
	b=1
	\]

	Así:

	\[
	B=
	\{x+cx^2:c\in\R\}
	\]

	El polinomio nulo no pertenece a $B$.

	\[
	\boxed{
	B\text{ no es subespacio vectorial}
	}
	\]

\section{Ejercicio 6}
	Sea:

	\[
	M=
	\left\{
	(x_1,x_2,x_3,x_4)\in\R^4:
	\begin{array}{rcl}
	x_1+5x_2-2x_3-x_4&=&0\\
	2x_1+x_2+2x_3+x_4&=&0
	\end{array}
	\right\}
	\]

	\subsection{Solución}

	Tomamos:

	\[
	x_3=s,
	\qquad
	x_4=t
	\]

	Resolviendo el sistema:

	\[
	x_1
	=
	-\frac{4}{3}s
	-\frac{2}{3}t
	\]

	\[
	x_2
	=
	\frac{2}{3}s
	+\frac{1}{3}t
	\]

	Por tanto:

	\[
	\boxed{
	\begin{aligned}
	x_1&=-\frac{4}{3}s-\frac{2}{3}t\\
	x_2&=\frac{2}{3}s+\frac{1}{3}t\\
	x_3&=s\\
	x_4&=t
	\end{aligned}
	\qquad s,t\in\R
	}
	\]

	Equivalentemente:

	\[
	M=
	\Span\left\{
	(-4,2,3,0),
	(-2,1,0,3)
	\right\}
	\]

\section{Ejercicio 7}
	Sean:

	\[
	S=
	\{(x_1,x_2,x_3,x_4)\in\R^4:x_1-x_2=0\}
	\]

	\[
	T=
	\Span\{(1,1,2,1),(2,3,-1,1)\}
	\]

	\subsection{Solución}

	Para $S$:

	\[
	x_1=x_2
	\]

	\[
	(x_1,x_2,x_3,x_4)
	=
	(x_1,x_1,x_3,x_4)
	\]

	Por tanto:

	\[
	\boxed{
	B_S=
	\{(1,1,0,0),(0,0,1,0),(0,0,0,1)\}
	}
	\]

	Los generadores de $T$ son linealmente independientes:

	\[
	\boxed{
	B_T=
	\{(1,1,2,1),(2,3,-1,1)\}
	}
	\]

	Un vector de $T$ es:

	\[
	(\alpha+2\beta,\alpha+3\beta,2\alpha-\beta,\alpha+\beta)
	\]

	Para pertenecer a $S$:

	\[
	\alpha+2\beta
	=
	\alpha+3\beta
	\]

	\[
	\beta=0
	\]

	Por tanto:

	\[
	S\cap T
	=
	\{\alpha(1,1,2,1):\alpha\in\R\}
	\]

	\[
	\boxed{
	B_{S\cap T}
	=
	\{(1,1,2,1)\}
	}
	\]

	Aplicamos Grassmann:

	\[
	\dimn(S+T)
	=
	3+2-1
	=
	4
	\]

	Así:

	\[
	S+T=\R^4
	\]

	y podemos tomar:

	\[
	\boxed{
	B_{S+T}
	=
	\{(1,0,0,0),(0,1,0,0),(0,0,1,0),(0,0,0,1)\}
	}
	\]

\section{Ejercicio 8}
	Analizar para qué valores de $a$ son linealmente independientes:

	\[
	A_1=
	\begin{pmatrix}
	0&a&0\\
	0&1&0
	\end{pmatrix},
	\quad
	A_2=
	\begin{pmatrix}
	0&1&0\\
	0&a&0
	\end{pmatrix},
	\quad
	A_3=
	\begin{pmatrix}
	0&1&0\\
	0&a&a
	\end{pmatrix}
	\]

	\subsection{Solución}

	Consideramos:

	\[
	\alpha A_1+\beta A_2+\gamma A_3=0
	\]

	De las entradas no nulas:

	\[
	\alpha a+\beta+\gamma=0
	\]

	\[
	\alpha+a\beta+a\gamma=0
	\]

	\[
	a\gamma=0
	\]

	Si $a=0$, el sistema es ligado.

	Si $a\neq0$:

	\[
	\gamma=0
	\]

	y:

	\[
	a\alpha+\beta=0
	\]

	\[
	\alpha+a\beta=0
	\]

	Sustituyendo:

	\[
	(1-a^2)\alpha=0
	\]

	Para que únicamente exista la solución trivial:

	\[
	a\neq\pm1
	\]

	Por tanto:

	\[
	\boxed{
	a\notin\{-1,0,1\}
	}
	\]

\section{Ejercicio 9}
	Analizar para qué valores de $a$ son linealmente independientes:

	\[
	x^2+3x+1,
	\qquad
	2-x,
	\qquad
	1+ax+x^2
	\]

	\subsection{Solución}

	En la base canónica de $\Pn_2(x)$:

	\[
	v_1=(1,3,1)
	\]

	\[
	v_2=(2,-1,0)
	\]

	\[
	v_3=(1,a,1)
	\]

	Calculamos:

	\[
	\det
	\begin{pmatrix}
	1&3&1\\
	2&-1&0\\
	1&a&1
	\end{pmatrix}
	=
	2(a-3)
	\]

	Los vectores son linealmente independientes si el determinante es distinto de cero:

	\[
	2(a-3)\neq0
	\]

	\[
	\boxed{
	a\neq3
	}
	\]

\section{Ejercicio 10}
	Sea:

	\[
	A=
	\begin{pmatrix}
	1&-1&1\\
	1&-1&1
	\end{pmatrix}
	\]

	y:

	\[
	U=
	\{X\in M_{2\times2}(\R):XA=0\}
	\]

	\subsection{Solución}

	Tomamos:

	\[
	X=
	\begin{pmatrix}
	x&y\\
	z&t
	\end{pmatrix}
	\]

	Entonces:

	\[
	XA=
	\begin{pmatrix}
	x+y&-(x+y)&x+y\\
	z+t&-(z+t)&z+t
	\end{pmatrix}
	\]

	La condición $XA=0$ implica:

	\[
	x+y=0
	\]

	\[
	z+t=0
	\]

	Por tanto:

	\[
	X=
	\begin{pmatrix}
	x&-x\\
	z&-z
	\end{pmatrix}
	\]

	\[
	=
	x
	\begin{pmatrix}
	1&-1\\
	0&0
	\end{pmatrix}
	+
	z
	\begin{pmatrix}
	0&0\\
	1&-1
	\end{pmatrix}
	\]

	Así:

	\[
	\boxed{\dimn U=2}
	\]

	\[
	\boxed{
	B_U=
	\left\{
	\begin{pmatrix}
	1&-1\\
	0&0
	\end{pmatrix},
	\begin{pmatrix}
	0&0\\
	1&-1
	\end{pmatrix}
	\right\}
	}
	\]

\section{Ejercicio 11}
	En $\R^4$ calcular las bases pedidas.

	\subsection{Solución}

	\textbf{a)}

	Partimos de:

	\[
	v=(1,2,1,1)
	\]

	Añadimos tres vectores canónicos:

	\[
	e_1=(1,0,0,0),
	\quad
	e_2=(0,1,0,0),
	\quad
	e_3=(0,0,1,0)
	\]

	Los cuatro vectores son linealmente independientes.

	\[
	\boxed{
	B=
	\{(1,2,1,1),(1,0,0,0),(0,1,0,0),(0,0,1,0)\}
	}
	\]

	\textbf{b)}

	Los tres vectores dados tienen rango $3$:

	\[
	\rg
	\begin{pmatrix}
	1&1&0&0\\
	0&0&2&2\\
	0&3&3&0
	\end{pmatrix}
	=
	3
	\]

	Añadimos, por ejemplo:

	\[
	(0,0,0,1)
	\]

	Por tanto:

	\[
	\boxed{
	B=
	\{(1,1,0,0),(0,0,2,2),(0,3,3,0),(0,0,0,1)\}
	}
	\]

\section{Ejercicio 12}
	Sean:

	\[
	B=\{e_1,e_2,e_3\}
	\]

	\[
	B'=
	\{2e_1+3e_2,\ e_1+e_3,\ -e_2+e_3\}
	\]

	\subsection{Solución}

	La matriz que transforma coordenadas de $B'$ a $B$ es:

	\[
	M_{B'B}
	=
	\begin{pmatrix}
	2&1&0\\
	3&0&-1\\
	0&1&1
	\end{pmatrix}
	\]

	Su inversa es:

	\[
	M_{BB'}
	=
	\begin{pmatrix}
	-1&1&1\\
	3&-2&-2\\
	-3&2&3
	\end{pmatrix}
	\]

	\textbf{a)}

	\[
	[u]_{B'}
	=
	M_{BB'}[u]_B
	\]

	\[
	=
	\begin{pmatrix}
	-1&1&1\\
	3&-2&-2\\
	-3&2&3
	\end{pmatrix}
	\begin{pmatrix}
	1\\2\\3
	\end{pmatrix}
	\]

	\[
	=
	\begin{pmatrix}
	4\\-7\\10
	\end{pmatrix}
	\]

	\[
	\boxed{
	[u]_{B'}=(4,-7,10)
	}
	\]

	\textbf{b)}

	\[
	[u]_B
	=
	M_{B'B}[u]_{B'}
	\]

	\[
	=
	\begin{pmatrix}
	2&1&0\\
	3&0&-1\\
	0&1&1
	\end{pmatrix}
	\begin{pmatrix}
	-2\\1\\0
	\end{pmatrix}
	\]

	\[
	=
	\begin{pmatrix}
	-3\\-6\\1
	\end{pmatrix}
	\]

	\[
	\boxed{
	[u]_B=(-3,-6,1)
	}
	\]

\section{Ejercicio 13}
	Sean las bases $B$, $B'$ y $B''$ del enunciado.

	\subsection{Solución}

	La matriz cuyas columnas son los vectores de $B'$ en la base canónica $B$ es:

	\[
	P=
	\begin{pmatrix}
	1&0&0&1\\
	1&1&0&0\\
	0&1&1&0\\
	0&0&-1&2
	\end{pmatrix}
	\]

	Por tanto:

	\[
	M_{BB'}
	=
	P^{-1}
	\]

	\[
	\boxed{
	M_{BB'}
	=
	\frac{1}{3}
	\begin{pmatrix}
	2&1&-1&-1\\
	-2&2&1&1\\
	2&-2&2&-1\\
	1&-1&1&1
	\end{pmatrix}
	}
	\]

	Para $B''$:

	\[
	Q=
	\begin{pmatrix}
	1&1&1&1\\
	0&1&1&1\\
	0&0&1&1\\
	0&0&0&1
	\end{pmatrix}
	\]

	Entonces:

	\[
	M_{BB''}
	=
	Q^{-1}
	\]

	\[
	\boxed{
	M_{BB''}
	=
	\begin{pmatrix}
	1&-1&0&0\\
	0&1&-1&0\\
	0&0&1&-1\\
	0&0&0&1
	\end{pmatrix}
	}
	\]

	Además:

	\[
	M_{B'B''}
	=
	M_{BB''}M_{B'B}
	\]

	\[
	\boxed{
	M_{B'B''}
	=
	\begin{pmatrix}
	0&-1&0&1\\
	1&0&-1&0\\
	0&1&2&-2\\
	0&0&-1&2
	\end{pmatrix}
	}
	\]

	Para:

	\[
	p(x)=x^3-2x
	\]

	se tiene:

	\[
	[p]_B=(0,-2,0,1)
	\]

	\[
	[p]_{B'}
	=
	M_{BB'}[p]_B
	=
	(-1,-1,1,1)
	\]

	\[
	[p]_{B''}
	=
	M_{BB''}[p]_B
	=
	(2,-2,-1,1)
	\]

	Por tanto:

	\[
	\boxed{
	[p]_B=(0,-2,0,1),
	\quad
	[p]_{B'}=(-1,-1,1,1),
	\quad
	[p]_{B''}=(2,-2,-1,1)
	}
	\]

\section{Ejercicio 14}
	Sean las bases $B_1$ y $B_2$ del enunciado.

	\subsection{Solución}

	\textbf{a)}

	En coordenadas canónicas:

	\[
	B_1=
	\{(1,1,1),(0,1,2),(1,1,0)\}
	\]

	\[
	B_2=
	\{(1,2,3),(\alpha,\alpha-1,-2),(2,2,0)\}
	\]

	La condición:

	\[
	(1,1,0)_{B_2}
	=
	(1,0,3)_{B_1}
	\]

	implica:

	\[
	(1,2,3)+(\alpha,\alpha-1,-2)
	=
	(1,1,1)+3(1,1,0)
	\]

	\[
	(1+\alpha,1+\alpha,1)
	=
	(4,4,1)
	\]

	\[
	\boxed{\alpha=3}
	\]

	\textbf{b)}

	Con $\alpha=3$, un vector tiene las mismas coordenadas $(a,b,c)$ en ambas bases si:

	\[
	a(1,2,3)+b(3,2,-2)+c(2,2,0)
	\]

	\[
	=
	a(1,1,1)+b(0,1,2)+c(1,1,0)
	\]

	De aquí:

	\[
	(a,b,c)
	=
	(2\lambda,\lambda,-3\lambda)
	\]

	Por tanto:

	\[
	W=
	\Span\{(2,1,-3)\}
	\]

	\[
	\boxed{
	\dimn W=1
	}
	\]

	y $W$ es un subespacio vectorial.

	\textbf{c)}

	Un posible subespacio suplementario de $W$ viene dado, en coordenadas respecto de $B_1$, por:

	\[
	W_s=
	\{(x,y,z)_{B_1}:y=x+z\}
	\]

	\[
	\boxed{
	W\oplus W_s=\Pn_2(x)
	}
	\]

\section{Ejercicio 15}
	Sean $S_1$ y $S_2$ los subespacios del enunciado.

	\subsection{Solución}

	Tomamos:

	\[
	p(x)=a+bx+cx^2+dx^3
	\]

	Para $S_1$:

	\[
	p(0)=0
	\quad\Longrightarrow\quad
	a=0
	\]

	Las tangentes en $x=1$ y $x=-1$ son paralelas si:

	\[
	p'(1)=p'(-1)
	\]

	Como:

	\[
	p'(x)=b+2cx+3dx^2
	\]

	se obtiene:

	\[
	b+2c+3d
	=
	b-2c+3d
	\]

	\[
	c=0
	\]

	Así:

	\[
	S_1=
	\{bx+dx^3:b,d\in\R\}
	\]

	\[
	\boxed{
	\dimn S_1=2,
	\qquad
	B_{S_1}=\{x,x^3\}
	}
	\]

	Para $S_2$:

	\[
	p(2)=0
	\]

	\[
	a+2b+4c+8d=0
	\]

	\[
	a=-2b-4c-8d
	\]

	Por tanto:

	\[
	S_2=
	\Span\{-2+x,-4+x^2,-8+x^3\}
	\]

	\[
	\boxed{
	\dimn S_2=3
	}
	\]

	Para la intersección, los polinomios son:

	\[
	p(x)=bx+dx^3
	\]

	y además:

	\[
	p(2)=2b+8d=0
	\]

	\[
	b=-4d
	\]

	Así:

	\[
	S_1\cap S_2
	=
	\Span\{x^3-4x\}
	\]

	\[
	\boxed{
	B_{S_1\cap S_2}
	=
	\{x^3-4x\}
	}
	\]

	Como:

	\[
	S_1\cap S_2\neq\{0\}
	\]

	\[
	\boxed{
	S_1+S_2\text{ no es suma directa}
	}
	\]

	Consideramos el vector generador:

	\[
	p(x)=x^3-4x
	\]

	En la base:

	\[
	B=
	\{1+x,x+x^2,x^2+x^3,x^3\}
	\]

	planteamos:

	\[
	(0,-4,0,1)
	=
	\alpha(1,1,0,0)
	+\beta(0,1,1,0)
	+\gamma(0,0,1,1)
	+\delta(0,0,0,1)
	\]

	Se obtiene:

	\[
	\alpha=0,
	\qquad
	\beta=-4,
	\qquad
	\gamma=4,
	\qquad
	\delta=-3
	\]

	\[
	\boxed{
	[p]_B=(0,-4,4,-3)
	}
	\]

	La matriz de paso de $B$ a la base canónica es:

	\[
	M_{BB_c}
	=
	\begin{pmatrix}
	1&0&0&0\\
	1&1&0&0\\
	0&1&1&0\\
	0&0&1&1
	\end{pmatrix}
	\]

	y:

	\[
	M_{BB_c}
	\begin{pmatrix}
	0\\-4\\4\\-3
	\end{pmatrix}
	=
	\begin{pmatrix}
	0\\-4\\0\\1
	\end{pmatrix}
	\]

	que corresponde a:

	\[
	\boxed{
	p(x)=x^3-4x
	}
	\]

\section{Ejercicio 16}
	Sea:

	\[
	p(x)=a+bx+cx^2,
	\qquad
	c\neq0
	\]

	\subsection{Solución}

	\textbf{a)}

	Se tiene:

	\[
	p'(x)=b+2cx
	\]

	\[
	p''(x)=2c
	\]

	En la base canónica:

	\[
	p=(a,b,c)
	\]

	\[
	p'=(b,2c,0)
	\]

	\[
	p''=(2c,0,0)
	\]

	La matriz formada por estos vectores tiene rango $3$ porque $c\neq0$.

	Por tanto:

	\[
	\boxed{
	B=\{p,p',p''\}\text{ es base de }\Pn_2(x)
	}
	\]

	Para:

	\[
	B'=\{p,p+p',p'+p''\}
	\]

	restando vectores se obtiene el mismo sistema generado por $\{p,p',p''\}$, por lo que también tiene rango $3$.

	\[
	\boxed{
	B'\text{ es base de }\Pn_2(x)
	}
	\]

	\textbf{b)}

	\[
	S=
	\{x(x-a):a\in\R\}
	\]

	\[
	=
	\{x^2-ax:a\in\R\}
	\]

	En coordenadas:

	\[
	S=
	\{(0,-a,1):a\in\R\}
	\]

	El vector nulo no pertenece a $S$.

	\[
	\boxed{
	S\text{ no es subespacio vectorial}
	}
	\]

	\textbf{c)}

	Si $p(x)\in S$ y $-1$ es raíz:

	\[
	p(x)=x^2-ax
	\]

	\[
	p(-1)=1+a=0
	\]

	\[
	a=-1
	\]

	Así:

	\[
	p(x)=x^2+x
	\]

	\[
	p'(x)=1+2x
	\]

	\[
	p''(x)=2
	\]

	La base $B'$ es:

	\[
	B'=
	\{x+x^2,\;1+3x+x^2,\;3+2x\}
	\]

	Buscamos:

	\[
	q(x)=x^2+x+2
	\]

	como:

	\[
	q=
	\alpha(x+x^2)
	+
	\beta(1+3x+x^2)
	+
	\gamma(3+2x)
	\]

	Se obtiene:

	\[
	\alpha=2,
	\qquad
	\beta=-1,
	\qquad
	\gamma=1
	\]

	\[
	\boxed{
	[q]_{B'}=(2,-1,1)
	}
	\]

	\textbf{d)}

	La matriz de cambio de $B'$ a la base canónica es:

	\[
	M_{B'B_c}
	=
	\begin{pmatrix}
	0&1&3\\
	1&3&2\\
	1&1&0
	\end{pmatrix}
	\]

	\[
	M_{B'B_c}
	\begin{pmatrix}
	2\\-1\\1
	\end{pmatrix}
	=
	\begin{pmatrix}
	2\\1\\1
	\end{pmatrix}
	\]

	Por tanto:

	\[
	\boxed{
	q(x)=2+x+x^2
	}
	\]

\section{Ejercicio 17}
	Sean:

	\[
	V=
	\Span\{(1,2,3,4),(-1,0,1,-1)\}
	\]

	\[
	W=
	\{(x,y,z,t)\in\R^4:2x+5y-z-t=0\}
	\]

	\subsection{Solución}

	\textbf{a)}

	Despejamos:

	\[
	t=2x+5y-z
	\]

	\[
	(x,y,z,t)
	=
	x(1,0,0,2)
	+y(0,1,0,5)
	+z(0,0,1,-1)
	\]

	\[
	\boxed{
	B_W=
	\{(1,0,0,2),(0,1,0,5),(0,0,1,-1)\}
	}
	\]

	\textbf{b)}

	Un vector de $V$ tiene la forma:

	\[
	(\alpha-\beta,2\alpha,3\alpha+\beta,4\alpha-\beta)
	\]

	Para pertenecer a $W$:

	\[
	2(\alpha-\beta)
	+
	5(2\alpha)
	-
	(3\alpha+\beta)
	-
	(4\alpha-\beta)
	=
	0
	\]

	\[
	5\alpha-4\beta=0
	\]

	\[
	\alpha=\frac45\beta
	\]

	Por tanto:

	\[
	V\cap W
	=
	\Span\{(-1,8,17,11)\}
	\]

	\[
	\boxed{
	B_{V\cap W}
	=
	\{(-1,8,17,11)\}
	}
	\]

	Una expresión paramétrica es:

	\[
	\boxed{
	(x,y,z,t)
	=
	\lambda(-1,8,17,11)
	}
	\]

	y una expresión implícita:

	\[
	\boxed{
	y+8x=0,
	\qquad
	z+17x=0,
	\qquad
	t+11x=0
	}
	\]

	\textbf{c)}

	Como:

	\[
	V\cap W\neq\{0\}
	\]

	\[
	\boxed{
	V+W\text{ no es suma directa}
	}
	\]

	\textbf{d)}

	Planteamos:

	\[
	(-1,1,1,2)
	=
	\alpha(1,2,3,4)
	+
	\beta(-1,0,1,-1)
	\]

	El sistema es incompatible.

	\[
	\boxed{
	(-1,1,1,2)\notin V
	}
	\]

	Por tanto, no existen coordenadas de dicho vector en una base de $V$.

\section{Ejercicio 18}
	Sean $V$ y $W$ los subespacios del enunciado.

	\subsection{Solución}

	Sea:

	\[
	p(x)=a+bx+cx^2+dx^3
	\]

	Entonces:

	\[
	p'(x)=b+2cx+3dx^2
	\]

	Para que:

	\[
	p'(x)\in
	\Span\{1+x^2,x^3\}
	\]

	como $p'(x)$ tiene grado menor o igual que $2$, debe cumplirse:

	\[
	p'(x)=\alpha(1+x^2)
	\]

	Por tanto:

	\[
	b=\alpha,
	\qquad
	c=0,
	\qquad
	3d=\alpha
	\]

	\[
	b=3d
	\]

	Así:

	\[
	V=
	\{(a,3d,0,d):a,d\in\R\}
	\]

	\[
	\boxed{
	B_V=
	\{(1,0,0,0),(0,3,0,1)\}
	}
	\]

	y sus ecuaciones implícitas son:

	\[
	\boxed{
	c=0,
	\qquad
	b-3d=0
	}
	\]

	Para $W$:

	\[
	p'(x)=p''(x)
	\]

	\[
	b+2cx+3dx^2
	=
	2c+6dx
	\]

	Comparando coeficientes:

	\[
	b=2c,
	\qquad
	2c=6d,
	\qquad
	3d=0
	\]

	\[
	b=c=d=0
	\]

	Por tanto:

	\[
	W=
	\{(a,0,0,0):a\in\R\}
	\]

	\[
	\boxed{
	B_W=\{(1,0,0,0)\}
	}
	\]

	Como:

	\[
	W\subset V
	\]

	se tiene:

	\[
	V\cap W=W
	\]

	\[
	\boxed{
	B_{V\cap W}
	=
	\{(1,0,0,0)\}
	}
	\]

	La expresión paramétrica es:

	\[
	\boxed{
	(a,0,0,0),
	\qquad
	a\in\R
	}
	\]

	El polinomio $p(x)=1$ pertenece a $V\cap W$, y por tanto también a $V$.

	En la base obtenida de $V$:

	\[
	\boxed{
	[1]_{B_V}=(1,0)
	}
	\]

	En la base de $V\cap W$:

	\[
	\boxed{
	[1]_{B_{V\cap W}}=(1)
	}
	\]

\section{Ejercicio 19}
	Sea $\Pn_n(x)$.

	\subsection{Solución}

	\textbf{a)}

	Consideramos:

	\[
	x^n,
	\quad
	nx^{n-1},
	\quad
	n(n-1)x^{n-2},
	\quad\ldots,\quad
	n!
	\]

	En la base canónica:

	\[
	\{1,x,\ldots,x^n\}
	\]

	cada uno tiene un único coeficiente principal no nulo en una posición distinta.

	Por tanto, son $n+1$ vectores linealmente independientes de un espacio de dimensión $n+1$.

	\[
	\boxed{
	\{x^n,(x^n)',\ldots,(x^n)^{(n)}\}
	\text{ es una base de }\Pn_n(x)
	}
	\]

	\textbf{b)}

	En coordenadas:

	\[
	v_1=(1,3,5)
	\]

	\[
	v_2=(-1,0,2)
	\]

	\[
	v_3=(3,3,1)
	\]

	Se tiene:

	\[
	\rg
	\begin{pmatrix}
	1&3&5\\
	-1&0&2\\
	3&3&1
	\end{pmatrix}
	=
	2
	\]

	Por tanto:

	\[
	\boxed{
	v_1,v_2,v_3
	\text{ son linealmente dependientes}
	}
	\]

	\textbf{c)}

	\[
	V=
	\Span\{(1,0,1),(1,0,-1)\}
	\]

	Para:

	\[
	p(x)=1+5x^2
	\]

	\[
	(1,0,5)
	=
	3(1,0,1)
	-2(1,0,-1)
	\]

	\[
	\boxed{
	p(x)\in V
	}
	\]

	Para:

	\[
	r(x)=1+x
	\]

	\[
	(1,1,0)
	\]

	no puede expresarse como combinación lineal de los generadores de $V$, ya que todos ellos tienen coeficiente de $x$ igual a cero.

	\[
	\boxed{
	r(x)\notin V
	}
	\]

	\textbf{d)}

	Como el tercer generador de $W$ depende de los dos primeros:

	\[
	\dimn W=2
	\]

	Podemos escribir:

	\[
	W=
	\Span\{(1,3,5),(-1,0,2)\}
	\]

	Un vector de $V$ tiene forma:

	\[
	(\alpha+\beta,0,\alpha-\beta)
	\]

	y uno de $W$:

	\[
	(\gamma-\delta,3\gamma,5\gamma+2\delta)
	\]

	Igualando:

	\[
	\gamma=0
	\]

	\[
	\delta=2\alpha
	\]

	\[
	\beta=-3\alpha
	\]

	Así:

	\[
	V\cap W
	=
	\Span\{(-2,0,4)\}
	\]

	\[
	\boxed{
	V\cap W
	=
	\Span\{-2+4x^2\}
	}
	\]

	Por Grassmann:

	\[
	\dimn(V+W)
	=
	2+2-1
	=
	3
	\]

	Por tanto:

	\[
	V+W=\Pn_2(x)
	\]

	Una base es:

	\[
	\boxed{
	B_{V+W}
	=
	\{1+x^2,\;1-x^2,\;1+3x+5x^2\}
	}
	\]

\section{Ejercicio 20}
	Sean $E=F=\R^2$.

	\subsection{Solución}

	La base canónica de $\R^2$ es:

	\[
	e_1=(1,0),
	\qquad
	e_2=(0,1)
	\]

	Una base de $E\times F$ es:

	\[
	\{(e_1,0),(e_2,0),(0,e_1),(0,e_2)\}
	\]

	En coordenadas:

	\[
	\boxed{
	B_{E\times F}
	=
	\{(1,0,0,0),(0,1,0,0),(0,0,1,0),(0,0,0,1)\}
	}
	\]

	\[
	\boxed{
	\dimn(E\times F)=4
	}
	\]

\section{Ejercicio 21}
	Sea $\dimn E=3$ y:

	\[
	V=
	\Span\{e_1-e_2,e_1+e_2\}
	\]

	\subsection{Solución}

	Los dos generadores de $V$ son linealmente independientes:

	\[
	\dimn V=2
	\]

	Por tanto:

	\[
	\dimn(E/V)
	=
	3-2
	=
	1
	\]

	Necesitamos una clase no nula de $E/V$, es decir, la clase de un vector que no pertenezca a $V$.

	Podemos tomar $e_3$:

	\[
	e_3\notin V
	\]

	Por tanto:

	\[
	\boxed{
	B_{E/V}
	=
	\{[e_3]\}
	}
	\]

\section{Ejercicio 22}
	Sea:

	\[
	V=
	\Span\{(1,0,1,1),(1,2,1,1),(2,2,2,2)\}
	\]

	\subsection{Solución}

	Calculamos:

	\[
	\rg
	\begin{pmatrix}
	1&0&1&1\\
	1&2&1&1\\
	2&2&2&2
	\end{pmatrix}
	=
	2
	\]

	Por tanto:

	\[
	\dimn V=2
	\]

	y:

	\[
	\dimn(\R^4/V)
	=
	4-2
	=
	2
	\]

	Además:

	\[
	V=
	\Span\{(1,0,1,1),(1,2,1,1)\}
	\]

	Un vector de $V$ tiene la forma:

	\[
	(\alpha+\beta,2\beta,\alpha+\beta,\alpha+\beta)
	\]

	por lo que:

	\[
	x=z=t
	\]

	Los vectores:

	\[
	(1,0,0,0)
	\]

	y:

	\[
	(0,0,0,1)
	\]

	no pertenecen a $V$ y sus clases son linealmente independientes módulo $V$.

	Por tanto:

	\[
	\boxed{
	B_{\R^4/V}
	=
	\left\{
	[(1,0,0,0)],
	[(0,0,0,1)]
	\right\}
	}
	\]

\end{document}
